\subsection{Framing Prescription and 3-Manifold Identification}

Let \(\beta \in B_n\) be the braid word on \(n\) strands constructed in the previous subsection from a length-\(n\) periodic orbit \(\gamma\) of the shift \(\sigma_p\) on the Bruhat--Tits tree \(T_p\). We work throughout in the ribbon category \(\mathcal{C}_p = \mathrm{SU}(2)_{p-2}\) at \(q = e^{\pi i / p}\), whose objects are the simple modules \(V_j\) for \(j = 0, 1/2, \dots, (p-2)/2\), with ribbon element \(\theta_j = \exp\bigl(\pi i (j(j+1) - 1/4)/p\bigr)\) (see Section~\ref{sec:mtc} for the precise normalization of the \(T\)-matrix).

\subsubsection{Canonical framing}

We equip the braid diagram of \(\beta\) with the \emph{blackboard framing}: each strand is framed by the normal vector field lying in the plane of the diagram, pointing to the right when traversing the strand in the positive direction. Let \(w(\beta) \in \mathbb{Z}\) denote the writhe of this diagram (the number of positive crossings minus the number of negative crossings). The blackboard framing induces a framing number \(f_{\mathrm{bb}}(\beta) = w(\beta)\) on the closure \(\widehat{\beta}\).

To obtain a well-defined Reshetikhin--Turaev invariant independent of the choice of diagram (up to the usual Kirby moves), we apply the \emph{writhe correction}: the canonical framed link is the pair \((\widehat{\beta}, f)\) where the framing \(f\) is the blackboard framing adjusted by \(-w(\beta)\), i.e.,
\[
f = f_{\mathrm{bb}}(\beta) - w(\beta).
\]
In the language of the ribbon category \(\mathcal{C}_p\), this corresponds to tensoring the braid representation \(\rho(\beta) \in \mathrm{End}(V_{j_1} \otimes \cdots \otimes V_{j_n})\) with the ribbon element raised to the power of the total framing on each component:
\[
\mathrm{RT}_{\mathcal{C}_p}(\beta, f) = \frac{1}{D_p^{n-1}} \sum_{j_1,\dots,j_n} \bigl( \theta_{j_1}^{f_1} \cdots \theta_{j_n}^{f_n} \bigr) \langle \rho(\beta) \rangle_{j_1 \dots j_n},
\]
where \(D_p = \sqrt{p / (2 \sin^2(\pi/p))}\) is the total quantum dimension, the sum runs over admissible colorings (i.e., those satisfying the fusion rules at each crossing), and \(f_k\) is the framing number on the \(k\)-th component after correction. By construction, the total framing vector \((f_1, \dots, f_n)\) satisfies \(\sum f_k = 0\); this is the \emph{zero-framed closure} relative to the blackboard diagram.

This prescription ensures invariance under Reidemeister moves of type~I (which alter the writhe by \(\pm 1\) and multiply by \(\theta_j^{\pm 1}\)) and type~II/III (preserved by the braiding). The resulting invariant \(\mathrm{RT}_{\mathcal{C}_p}(\beta, f)\) is therefore well-defined on the isotopy class of the framed link.

\subsubsection{Identification as a Seifert fibered 3-manifold}

We now prove that the closure \(\widehat{\beta}\) (with the canonical framing \(f\)) is isotopic to a Seifert fibered 3-manifold. The key geometric fact is that the periodic orbit \(\gamma\) of \(\sigma_p\) on \(T_p\) projects, via the boundary map \(\partial T_p \to \mathbb{P}^1(\mathbb{Q}_p)\), to a cyclic permutation of \(n\) distinguished points on the projective line. More precisely, choose a base vertex \(v_0 \in V(T_p)\) and label the \(n\) geodesic rays emanating from the orbit points by strands \(1,\dots,n\) in the natural cyclic order induced by the action of \(\sigma_p\) (which multiplies the local coordinate by \(p\) modulo the lattice). The braid word \(\beta\) is then conjugate in \(B_n\) to the power of a full twist on the \(n\) strands, twisted by the local action of the generator of the stabilizer in \(\mathrm{PGL}(2,\mathbb{Q}_p)\).

By a standard argument in geometric topology (cf.\ the classification of Seifert fibered spaces via their fundamental groups and the fact that the monodromy is a finite-order rotation in the base orbifold), the complement of \(\widehat{\beta}\) in \(S^3\) admits a Seifert fibration over the base orbifold \(S^2\) with four exceptional fibers. Explicitly:

- The base is the 4-punctured sphere \(S^2 \setminus \{p_1, p_2, p_3, p_4\}\), where the punctures correspond to the two fixed points at infinity (the ends of the tree) and the two branch points arising from the projection of the orbit under the action of \(\sigma_p\).
- The regular fiber is the circle \(S^1\) obtained by flowing along the geodesic rays in \(T_p\).
- The exceptional fibers arise from the local stabilizers: two of multiplicity \(p\) (corresponding to the \(p\)-adic ramification at the ends) and two of multiplicity \(n\) (corresponding to the period of the orbit).

The Seifert invariants are therefore
\[
M = M\bigl(0; (p,1), (p,1), (n, b_1), (n, b_2)\bigr),
\]
where the Euler number \(e(M)\) of the fibration is
\[
e(M) = -\frac{2}{p} - \frac{b_1 + b_2}{n} + \frac{2}{n},
\]
and the integers \(b_1, b_2 \in \{0,1\}\) are determined by the parity of the writhe correction and the orientation of the cyclic permutation (explicitly, \(b_1 = w(\beta) \bmod n\), \(b_2 = 1 - b_1\) after the writhe correction ensures the total twisting is normalized). The normalized Euler number satisfies \(|e(M)| = 1/p + O(1/n)\), but the precise value is not needed for the RT computation; only the multiplicities enter the Verlinde formula.

To establish the identification rigorously, note that the fundamental group \(\pi_1(M)\) is generated by the exceptional fibers with relations coming from the central extension by the regular fiber, which matches exactly the presentation obtained from the braid group closure under the cyclic monodromy. Isotopy invariance follows from the fact that any two such diagrams (differing by a choice of base vertex) differ by a sequence of Reidemeister moves and Kirby moves that preserve the Seifert structure.

\subsubsection{Independence under choice of distinguished end}

Different choices of the distinguished end at \(\infty \in \mathbb{P}^1(\mathbb{Q}_p)\) (i.e., different fixed points for the shift \(\sigma_p\)) yield Galois-conjugate Seifert manifolds. More precisely, the action of the absolute Galois group \(\mathrm{Gal}(\overline{\mathbb{Q}}_p / \mathbb{Q}_p)\) on the boundary permutes the possible ends, and the resulting braid words \(\beta'\) are related to \(\beta\) by an outer automorphism of \(B_n\) that preserves the conjugacy class up to Galois action on the cyclotomic field \(\mathbb{Q}(\zeta_p)\). Consequently, the Reshetikhin--Turaev invariants satisfy
\[
\mathrm{RT}_{\mathcal{C}_p}(\beta', f') = \sigma\bigl( \mathrm{RT}_{\mathcal{C}_p}(\beta, f) \bigr)
\]
for some \(\sigma \in \mathrm{Gal}(\mathbb{Q}(\zeta_p)/\mathbb{Q})\). In particular,
\[
\bigl| \mathrm{RT}_{\mathcal{C}_p}(\beta, f) \bigr| = \bigl| \mathrm{RT}_{\mathcal{C}_p}(\beta', f') \bigr|,
\]
so the absolute value (which is all that enters the local transfer operator trace) is independent of the auxiliary choice of end. This completes the framing and manifold identification, preparing the ground for the Verlinde reduction in Section~\ref{sec:proof}.

\begin{remark}
The explicit Seifert invariants above are independent of the particular representative of the orbit \(\gamma\) (all length-\(n\) orbits are conjugate under the automorphism group of \(T_p\)). Numerical checks for small \(p\) and \(n\) (see Appendix~C) confirm that the computed Euler number and multiplicities reproduce the expected multiplicities in the Verlinde sum.
\end{remark}